% Vorspann der Arbeit
%

\includegraphics[width=45mm]{pictures/Hochschule_Muenchen_Logo.png}
\hfill
\includegraphics[width=45mm]{pictures/ARRI_AG_Corporate_Logo.png}

\vspace*{20mm}
\begin{center}
{\large Hochschule f\"ur angewandte Wissenschaften M\"unchen}\\
{\large Fakult\"at f\"ur Elektrotechnik und Informationstechnik}\\
{\large Bachelorstudiengang Elektrotechnik und Informationstechnik}\\

\vspace*{15mm}
{\huge Bachelorarbeit }    %% oder Masterarbeit
\\

\vspace*{10mm}
{\huge \bfseries{%
                      \par Modulare Ansteuerung eines FPGA \"uber Software \par
}} 


\vspace*{15mm}
{\Large abgegeben von %
                      Maren Konrad 
%%                    Vorname Name
} 
\\
\end{center}

\vfill
{\large
\begin{tabbing}
\\
Bearbeitungsbeginn: \hspace{3.5cm} \= %
					  12.08.2019
\\
Abgabetermin: \> \= %     
                      20.01.2020
%%                    durch Ende-Datum ersetzen
\\
lfd. Nr. gem\"a{\ss}  Belegschein: \> %
					  1883 
\\
\end{tabbing} }
\thispagestyle{empty}
\cleardoublepage


% Seite 2
\vspace*{20mm}
\begin{center}
{\large Hochschule f\"ur angewandte Wissenschaften M\"unchen}\\
{\large Fakult\"at f\"ur Elektrotechnik und Informationstechnik}\\
{\large Bachelorstudiengang Elektrotechnik und Informationstechnik}\\

\vspace*{15mm}
{\huge Bachelorarbeit }    %% oder Masterarbeit
\\

\vspace*{10mm}
{\huge \bfseries{%
		\par Modulare Ansteuerung eines bildverarbeitenden FPGA "uber generische Kernelmodule \par
}} % deuscher Langtitel

\vspace*{10mm}
{\huge \bfseries{%
		\par Modular control of an image processing FPGA via generic kernel modules \par
}} % englischer Langtitel

\vspace*{15mm}
{\Large abgegeben von %
	Maren Konrad 
	%%                    Vorname Name
} 
\\
\end{center}

\vfill
{\large
	\begin{tabbing}
		\\
		Bearbeitungsbeginn: \hspace{3.5cm} \= %     
		12.08.2019
		%%                    durch Start-Datum ersetzen
		\\
		Abgabetermin: \> %     
		20.01.2020
		%%                    durch Ende-Datum ersetzen
		\\
		lfd. Nr. gem\"a{\ss}  Belegschein: \> %
		1883
		%%                    durch Lfd. Nr. ersetzen
		\\
		Betreuer (Hochschule M\"unchen): \> %         
		Prof. Dr. Gerhard Schillhuber
		%%                    hier den Prof. eintragen
		\\
		Betreuer (Extern): \> %         
		Anton Hattendorf
\end{tabbing} }
\thispagestyle{empty}
\cleardoublepage

% Seite 3
\thispagestyle{empty}
Erkl\"arungen des Bearbeiters:\\
\\
Name: Konrad \\
Vorname: Maren \\

1) Ich erkl\"are hiermit, dass ich die vorliegende Bachelorarbeit selbst\"andig verfasst
und noch nicht anderweitig zu Pr\"ufungszwecken vorgelegt habe.

S\"amtliche benutzte Quellen und Hilfsmittel sind angegeben, w\"ortliche und sinngem\"a{\ss}e Zitate sind als solche gekennzeichnet.
\vspace{15mm}

\begin{tabbing}
	M\"unchen, den \today \hspace{10mm} \= \rule{60mm}{0.5pt}\\
	\> {\small Unterschrift}
\end{tabbing}
\vspace{15mm}

\vfill
2) Der Ver\"offentlichung der Bachelorarbeit stimme ich hiermit \textbf{NICHT} zu.
\vspace{15mm}

\begin{tabbing}
	M\"unchen, den \today \hspace{10mm} \= \rule{60mm}{0.5pt}\\
	\> {\small Unterschrift}
\end{tabbing}

\vfill
\cleardoublepage

%Seite 4

\begin{center}
{\Large \bfseries{ Kurzfassung }}\\
\end{center}

%% Hier die Kurzfassung in Deutsch einf�gen
%In dem vorliegenden Bericht geht es um die Verbesserung einer Kamerasoftware. Es wird hierzu eine Abstraktionsebene vorgestellt, die aus mehreren Modulen besteht. Um das Verst\"andnis zu erleichtern wird au\ss{}erdem eine kurze Einf\"uhrung zur Bildkette gegeben und anschlie\ss{}end auf ausgew\"ahlte Module eingegangen. Durch die Abstraktionsebene wird die \"Ubersichtlichkeit und Erweiterung der Software in Zukunft erleichtert werden.


In dem vorliegenden Bericht geht es um die Verbesserung einer Kamerasoftware im Zusammenspiel mit der Hardware. Es wird hierzu eine Abstraktionsebene eingef\"uhrt, welche die Zugriffe von der Software auf die Module des FPGA \"ubersichtlicher und einfacher gestalten soll. Der Zugriff auf die Hardware findet aktuell direkt \"uber die Adresse des Moduls statt und soll abstrahiert werden, sodass die Software lediglich den Modulnamen ben\"otigt, um auf das richtige Modul zuzugreifen. Zum Erleichtern des Verst\"andnisses wird ein kurzer \"Uberblick \"uber die aktuelle Implementierung und eine kurze Einf\"uhrung zur Bildkette gegeben. Au\ss{}erdem wird auf das Betriebssystem Linux und ben\"otigte Komponenten eingegangen. Im Anschluss wird die Implementierung der Abstraktionsebene n\"aher erkl\"art. Die Wartbarkeit und die Erweiterung der Software sollen durch die Abstraktionsebene in Zukunft erleichtert werden.


Dieser Bericht richtet sich an Leser aus dem Bereich der Elektrotechnik.

\vspace*{5mm}
\begin{center}
{\Large \bfseries{ Abstract }}\\
\end{center}

%% Hier die Kurzfassung in Englisch ein
%This report handle the improvement of an camera software. For this, an abstraction level will be introduce, which consists out of a few modules. At first there is a short intorduction to the image processing chain to make the understanding easier and then go into details of selected modules. Because of the abstraction level the clarity and extensions of the software will be easier in future. 

This report decribes the improvement of a camera software in interaction with the hardware. For this purpose an abstraction layer will be introduced, which allows access from the software to the modules of the FPGA in a clearer and simpler way. The access to the hardware currently takes place directly via the address of the modules. This should be abstracted, so that the software only needs the name of a module to access the correct module. To facilitate understanding, a brief overview of the current implementation and a short introduction to the image chain are given. Besides, the basics of the Linux operating system and the required components are described. The implementaion of the abstraction layer is then explained in more detail. In the duture the maintainability and the extension of the software shall be facilitated by the abstraction layer.


This report targets readers having an electric engineering background.

\pagenumbering{Roman} 	
\renewcommand{\thechapter}{\Roman{chapter}} 

\cleardoublepage
%% Inhaltsverzeichnis
\tableofcontents
\clearpage



